\chapter{Finite-Temperature Path Integrals}
\label{ch:finite-temperature-path-integrals}

\ConventionsEuclidean

The KMS condition is the real-time equilibrium condition for a state and a
time evolution.  A Euclidean finite-temperature path integral is a regulated
coordinate representation of Gibbs-trace correlation functions when the trace
or a controlled infinite-volume replacement exists.  The trace derivation
explains the thermal circle and the boundary conditions; analytic
continuation back to real time remains a separate reconstruction problem.

\section{Trace to Euclidean Circle}

Let a quantum system have Hamiltonian \(H\) with
\(\operatorname{Tr}e^{-\beta H}<\infty\).  The thermal expectation is
\[
  \omega_\beta(A)
  =
  \frac{1}{Z(\beta)}\operatorname{Tr}(e^{-\beta H}A),
  \qquad
  Z(\beta)=\operatorname{Tr}e^{-\beta H}.
\]
Split \(\beta=N\epsilon\).  Inserting \(N\) resolutions of identity in a
configuration basis gives the Euclidean time-sliced expression
\[
  Z(\beta)
  =
  \lim_{N\to\infty}
  \int \prod_{j=1}^{N}dq_j\,
  \prod_{j=1}^{N}
  K_\epsilon(q_{j+1},q_j),
  \qquad
  q_{N+1}=q_1,
\]
where \(K_\epsilon(q',q)=\langle q'|e^{-\epsilon H}|q\rangle\).  The trace
imposes periodicity around the Euclidean time circle \(S^1_\beta\).

For a bosonic scalar field this becomes, after the spatial regulator and
continuum limit are specified,
\[
  Z(\beta)
  =
  \int_{\phi(\tau+\beta,\mathbf x)=\phi(\tau,\mathbf x)}
  [D\phi]_{\rm reg}\,
  e^{-S_E[\phi]} .
\]
The symbol \([D\phi]_{\rm reg}\) denotes the chosen regulated Euclidean
construction; the continuum meaning is supplied by the same constructive,
lattice, or perturbative framework used at zero temperature.

\section{Fermions and Spin Structure}

For fermions, the trace over a fermionic Fock space has a coherent-state
coordinate expression in Grassmann variables.  The cyclic trace produces
antiperiodic boundary conditions for the ordinary thermal partition function:
\[
  \psi(\tau+\beta,\mathbf x)=-\psi(\tau,\mathbf x),
  \qquad
  \bar\psi(\tau+\beta,\mathbf x)=-\bar\psi(\tau,\mathbf x).
\]
Equivalently, the thermal circle carries the antiperiodic spin structure.
The graded trace
\[
  \operatorname{Tr}\!\left((-1)^F e^{-\beta H}\cdots\right)
\]
uses the periodic spin structure.  This distinction is part of the definition
of the thermal observable and is essential in supersymmetric indices.

\section{Matsubara Frequencies}

On \(S^1_\beta\times\R^{D-1}\), periodic bosons have frequencies
\[
  \omega_n^{\rm B}=\frac{2\pi n}{\beta},
  \qquad n\in\mathbb Z,
\]
while antiperiodic fermions have
\[
  \omega_n^{\rm F}=\frac{(2n+1)\pi}{\beta}.
\]
For a free scalar,
\[
  G_E(\tau,\mathbf x)
  =
  \frac{1}{\beta}\sum_{n\in\mathbb Z}
  \int \frac{d^{D-1}k}{(2\pi)^{D-1}}\,
  \frac{e^{\ii\omega_n^{\rm B}\tau+\ii\mathbf k\cdot\mathbf x}}
       {(\omega_n^{\rm B})^2+\mathbf k^2+m^2}.
\]
Analytic continuation from discrete Euclidean frequencies to retarded
functions requires the spectral representation and the KMS condition.  The
Euclidean values at Matsubara frequencies alone become real-time response
data only after a growth condition or spectral ansatz fixes the continuation.

\section{Chemical Potentials and Background Holonomy}

For a conserved charge \(Q\), the grand-canonical trace is
\[
  Z(\beta,\mu)=\operatorname{Tr}e^{-\beta(H-\mu Q)} .
\]
In the Euclidean coordinate representation, \(\mu\) enters as an imaginary
background gauge-field holonomy along \(S^1_\beta\):
\[
  D_\tau\mapsto D_\tau-\mu Q .
\]
For compact gauge symmetries, a thermal holonomy is a gauge-invariant
variable modulo large gauge transformations, and the integration over gauge
holonomies is part of the gauge-theory finite-temperature path integral.

\begin{openproblem}[Euclidean thermal reconstruction]
\label{op:euclidean-thermal-reconstruction}
State constructive hypotheses under which Schwinger functions on
\(S^1_\beta\times\R^{D-1}\), with the appropriate spin structure and
background holonomies, reconstruct a KMS state with local real-time
observables and transport correlators.  The solution must specify the
analytic growth replacing the vacuum spectrum condition.
\end{openproblem}
